\documentclass[a4paper,11pt]{article}
\usepackage[utf8]{inputenc}
\usepackage[T1]{fontenc}
\usepackage[dutch]{babel}
\usepackage{geometry}
\geometry{margin=2.5cm}
\usepackage{booktabs}
\usepackage{tabularx}
\usepackage{xcolor}
\usepackage{titlesec}
\usepackage{microtype}
\usepackage{hyperref}

\definecolor{primary}{RGB}{28, 54, 84}
\definecolor{secondary}{RGB}{70, 95, 125}
\definecolor{accent}{RGB}{180, 80, 40}
\definecolor{lightgray}{gray}{0.95}

\titleformat{\section}{\Large\bfseries\color{primary}}{\thesection}{1em}{}[\titlerule]
\titleformat{\subsection}{\large\bfseries\color{secondary}}{\thesubsection}{1em}{}

\title{\textbf{\LARGE TRIFECTA Annotatie Pipeline}\\[0.3em] \Large Tussentijdse Analyse KWIC Batch}
\author{\textbf{Datum:} 28 augustus 2026 \quad $\cdot$ \quad \textbf{Model:} \texttt{qwen2.5-coder:latest}}
\date{}

\begin{document}

\maketitle

\vspace{-1.5em}
\noindent\rule{\textwidth}{1.5pt}
\vspace{1em}

\section{Overzicht \& Status}

Deze tussenrapportage vat de voortgang en tussenresultaten samen van de grootschalige annotatierun op het historische KWIC-snippetcorpus (\texttt{food\_snippets\_long\_kwic.csv}). 

\begin{itemize}
    \item \textbf{Totaal verwerkte records op peilmoment:} 9.148 records.
    \item \textbf{Executie-architectuur:} 3-traps LLM-pipeline (Stap A $\rightarrow$ Stap B $\rightarrow$ Stap C) met gestructureerde Pydantic/instructor schema's via Ollama (100\% GPU Metal-offload op Apple Silicon M4).
    \item \textbf{Vroegtijdige uitval (Dropout):} 2.995 records (32,7\%) zijn direct bij Stap A afgevallen, waardoor duizenden zware Stap B/C aanroepen zijn bespaard.
\end{itemize}

\section{Stap A: Entiteit- \& Metafoorfilter (Dropout)}

Stap A valideert of de doelterm in de historische context daadwerkelijk naar een voedingsentiteit verwijst en sluit metaforisch of oneigenlijk taalgebruik direct uit.

\begin{table}[h!]
\centering
\begin{tabular}{lrrp{7.5cm}}
\toprule
\textbf{Stap A Oordeel} & \textbf{Aantal} & \textbf{Percentage} & \textbf{Toelichting} \\
\midrule
\texttt{FOOD} (Bevestigd) & 6.153 & 67,3\% & Voedselcontext bevestigd; doorgestroomd naar Stap B \& C. \\
\texttt{DROPOUT} (Uitval) & 2.995 & 32,7\% & Geen voedselcontext, homoniem of metaforisch gebruik; vroege exit. \\
\midrule
\textbf{Totaal} & \textbf{9.148} & \textbf{100,0\%} & \\
\bottomrule
\end{tabular}
\caption{Verdeling van Stap A validatie.}
\end{table}

\section{Stap B: Macro-Frame Distributie}

Records die Stap A passeren worden in Stap B geclassificeerd in macro-frames. Onderstaand overzicht toont de relatieve aandelen over het gehele verwerkte corpus:

\begin{table}[h!]
\centering
\begin{tabular}{lrrp{6.5cm}}
\toprule
\textbf{Macro-Frame} & \textbf{Aantal} & \textbf{Aandeel} & \textbf{Karakteristiek / Typische Context} \\
\midrule
\texttt{COOKING\_CREATION} & 3.092 & 33,8\% & Recepten, kookbereidingen, culinaire handelingen. \\
\texttt{NONE} / \texttt{DROPOUT} & 2.995 & 32,7\% & Uitval via Stap A of niet-culinair domein. \\
\texttt{UNSPECIFIED} & 1.217 & 13,3\% & Voedselentiteit aanwezig zonder expliciete culinaire actie. \\
\texttt{CURE} & 852 & 9,3\% & Dieet- en medicinale receptuur, huismiddelen, remedies. \\
\texttt{INGESTION} & 787 & 8,6\% & Consumptie, eten, drinken, maaltijdsituaties. \\
\texttt{PRESERVING} & 205 & 2,2\% & Inmaken, pekelen, drogen, conserveren. \\
\midrule
\textbf{Totaal} & \textbf{9.148} & \textbf{100,0\%} & \\
\bottomrule
\end{tabular}
\caption{Distributie van Stap B macro-frames.}
\end{table}

\newpage

\section{Stap C: Qualia Slot-Invulling}

Voor de frames \texttt{COOKING\_CREATION}, \texttt{CURE}, \texttt{INGESTION} en \texttt{PRESERVING} worden in Stap C frame-specifieke Qualia-rollen geëxtraheerd.

\begin{table}[h!]
\centering
\begin{tabular}{llr}
\toprule
\textbf{Frame} & \textbf{Qualia Slot} & \textbf{Ingevulde Instanties} \\
\midrule
\textbf{COOKING\_CREATION} & \texttt{Method} (kookmethode / techniek) & 1.972 \\
& \texttt{Food\_Product} (bereid eindproduct) & 1.788 \\
& \texttt{Process} (procesverloop) & 1.360 \\
\midrule
\textbf{CURE} & \texttt{Food\_Treatment} (behandeling / toediening) & 784 \\
& \texttt{Affliction} (kwaal, aandoening, doel) & 671 \\
\midrule
\textbf{INGESTION} & \texttt{Manner} (wijze van consumptie) & 652 \\
& \texttt{Context} (maaltijd/setting) & 556 \\
& \texttt{Ingestor} (consument/persoon) & 489 \\
& \texttt{Food\_Patient} (geconsumeerd item) & 384 \\
& \texttt{Purpose} (doel van consumptie) & 249 \\
\midrule
\textbf{PRESERVING} & \texttt{PR\_Food\_Patient} (te conserveren waar) & 190 \\
& \texttt{PR\_Technique} (conserveringstechniek) & 152 \\
& \texttt{PR\_Medium} (inmaakmedium: azijn, pekel, suiker) & 72 \\
\bottomrule
\end{tabular}
\caption{Top ingevulde Qualia rollen per frame in Stap C.}
\end{table}

\section{Belangrijkste Inzichten \& Aanbevelingen}

\begin{enumerate}
    \item \textbf{Hoge densiteit in \texttt{COOKING\_CREATION} en \texttt{CURE}:} Zowel culinaire methoden als medicinale kwalen (\texttt{Affliction}) vertonen een zeer hoge slot-extractiegraad. Dit maakt de dataset direct bruikbaar voor diachrone analyses naar vroegmoderne geneeskundige voedingsrecepten vs. culinaire bereidingen.
    \item \textbf{Doeltreffendheid van de Dropout-architectuur:} Met een uitvalpercentage van 32,7\% bewijst de tweetraps filtering zijn waarde; 1 op de 3 records hoeft niet door de zware Qualia-prompting.
    \item \textbf{Schaalstrategie naar de toekomst:}
    \begin{itemize}
        \item \emph{Ontkoppelde Stap A filtering:} Voor nieuwe tranches eerst een bulk Stap A run draaien (\texttt{make step-a-bulk}).
        \item \emph{GijsBERT pre-filter:} Inzet van GijsBERT voor lokale high-throughput screening van \texttt{NONE}-gevallen.
        \item \emph{HPC Scaling (SURF):} Voor de volledige corpusverwerking ($>$100k snippets) opschalen naar dual NVIDIA A10 GPU's via \texttt{vLLM}.
    \end{itemize}
\end{enumerate}

\end{document}
